\chapter*{前言}\addcontentsline{toc}{chapter}{前言}\markboth{前言}{前言}
% 原书前言第1页

1708年12月4日，康熙皇帝一从北京城南郊的南苑回到宫中，太监立马通知在内廷养心殿附近制药、管理乐器的传教士早点离开，被废的太子允礽要奉召与康熙会面，地点就在养心殿。这是太子在10月17日宣布被废并遭圈禁后，父子两人第一次见面。傍晚时分，废太子被带到养心殿外院子中等候。根据传教士得知的消息，康熙父子二人的见面是“戏剧性”的：“只有泪水和抽泣”，都说不出话来。紧接着喇嘛来到了现场，开始做法事，为二人驱魔。“魔”先被驱赶到康熙和太子的衣服上，然后再从衣服上驱赶到驴的身上，最后通过鞭打驴来把魔驱赶出去。\NoteRef{preface}{1}


% 原书前言第1–2页

这一天是康熙四十七年十月二十三日乙丑。《圣祖仁皇帝实录》（即康熙朝《实录》，后简称《实录》）也记录了这一天康熙召见废太子，并写到康熙自陈召见废太子后心情的变化：“胸中亦不更有郁结矣。”《实录》中没有记载驱魔的法事。传教士在记录中也预感到关于法事的记录会让读者看起来“觉得可笑”，但驱“魔”在传教士眼中的确是祛除了康熙所谓的胸中“郁结”，这是太子从10月被废到下一年4月复立过程中的转折点。


% 原书前言第2页

从10月17日康熙宣布废黜太子算起，到12月4日父子相见的这四十八天中，夺嫡的斗争已经陆续上演：有见于传统清史记载、各种影视剧都演绎过的看相人张明德案；该案最后以张明德被凌迟处死、皇八子被削去贝勒爵位结案。还有皇三子举报大皇子派喇嘛用巫术镇魔太子的案子；最后以镇魔所用人偶自京中多处地下被挖出、大皇子被削去爵位和所有职务而结束。在传教士的记载中，这两个案子也是四十八天里的重头戏，但他们对两个案子中人物关系的记录与正统清史有很大差别。这是本书的重点，将会在后面仔细分析。更重要的是传教士还记录了正统清史中没有记录的部分：比如康熙身边的太监群体用计扳倒大皇子，使其失去了储位争夺的资格；还有皇四子，也就是后来的雍正皇帝，在这段时期的所作所为。


% 原书前言第2–3页

以《实录》为代表的官方记载在储位之争相关部分有所遗漏，以及记录之不可靠，是民国时期第一代清史研究者就已经发现的。最早系统论证清朝官方记录可靠性问题的是孟森。他是民国时期清史研究的奠基人，也是最早研究储位之争的学者。孟森在《清世宗入承大统考实》一文中，就通过比对《大义觉迷录》与《实录》，呈现出后者记载中的问题。\NoteRef{preface}{2}《大义觉迷录》在雍正七年（1729）由雍正亲自下令编纂，用以批判当时关于他得位不正的各种流言。雍正亲自在全国推广这本书。朝廷不仅颁发了官方版本，还要求各地方翻刻、推广，并安排专门人员在地方宣讲。雍正定下的宣传目标是：“各州县远乡僻壤”，“读书士子及乡曲小民”都要知晓。如果这个目标没有达到，雍正警告“定将该省学政及该县教官从重治罪”。在这样开动国家机器从上至下的宣传下，《大义觉迷录》成为了家喻户晓的书，在1720年代末到1730年代初广为流传。但刊布此书并没有达到雍正想要平息世间流言的目的，反而让关于他得位不正的传言流布更广。乾隆登基以后，在1736年全面禁毁此书。这本书流传到了日本，后来在清末被发现并重新印刷出版。\NoteRef{preface}{3}从《大义觉迷录》中，孟森发现了许多和《实录》有出入的记载。他提出了一系列萦绕在储位之争上的疑问：要是雍正的继位没有问题，那雍正朝的史官们为什么要在历史记载中做手脚？他们要达到什么目的？同时，这些问题又引发了对清廷官方记载可靠性的怀疑，到底清朝留下来的官方记录在储位之争相关历史上有多大程度的后期虚构？到底官方历史中哪些内容我们可以相信？


% 原书前言第3–4页

这些问题，在孟森以后百年来一直存在各种争议。\NoteRef{preface}{4}很重要的一个原因是与储位之争及太子废立相关的清宫原始档案极少。尽管近半个世纪以来，两岸学术机构整理出了大批清朝满汉档案，为我们全面了解清史提供了丰富的材料，但可惜的是，与储位之争相关的材料却并没有大量出现。比如本书提到的废黜太子后四十八天中发生的几起大案，留存下来的满汉奏折就没有直接相关的内容。同时，每天记载康熙日常事务的康熙朝《起居注》，偏偏这一年的没有存留。\NoteRef{preface}{5}而对于废太子问题，反而一直有传言《红楼梦》这部小说与之有关。小说开篇就以甄士隐、贾雨村两个虚构出的人物姓名的谐音点出：真事隐、假语存，来暗示世事记载有问题。更与本书有关的是，曹雪芹告诉我们：甄士隐，名费，字士隐，也就是甄费字士隐。现存最早的脂砚斋评本（甲戌本）中，“费”字旁边写了个“废”字，这样一来，谐音就变成了：真废之事隐。废太子事中是否真的别有隐情？我们看到的官方记述中有没有虚构情节？这些问题一直只有猜测，没有答案。


% 原书前言第4–5页

只有找到其他记录，然后再与清廷官方记载进行比较，才可能知道官方记载是否有虚构，以及虚构的程度。就像我们看到一张人像照片后，想要知道照片到底有多少是“美化”，以及“美化”的程度是否已经脱离了现实而成了虚构，那就只有找到其他拍摄者拍摄的照片才有可能破解。找寻官方记载以外的资料其实相当于在寻找其他的拍摄者。这是民国以来清史研究者一直在做的工作。大家最初想到的是在朝廷大臣的私人记录或者是从文人笔记中发现材料。比如吴秀良（Silas H. L. Wu）在1970年代出版的关于储位之争的英文专著中，提到《李文贞公（李光地）年谱》记载了康熙在外出行营路上决定废黜太子的导火索，是废太子有夜晚窥探康熙之举，所谓“帐殿夜警”让康熙感觉到了危险。\NoteRef{preface}{6}这条记录佐证了康熙朝《实录》记载的相关细节。


% 原书前言第5–6页

但总体来说，现在能找到的相关文人记录很少。即便找到，也都是像《李文贞公年谱》这样的只言片语。尽管当时废黜与复立太子都曾昭告天下，并非秘闻，但天下文人似乎集体噤声，各种文集笔记都鲜有记载。这可能与清代的文字狱有关。就算在康熙朝并非秘闻的废立太子之事，到了雍正和乾隆朝又成了禁忌。康熙朝著名诗人朱彝尊在他去世前几个月得闻太子复立，写了一首题为《三月十日诏下青宫再建喜而赋诗》的五言律诗。青色五行属木，而木为东方。青宫传统中指代太子的东宫。宣布太子复立在康熙四十八年三月初九日，而相关祭祀典礼在三月初十日。朱彝尊的诗揭示了复立太子是当时世人及时知晓的事件。这首诗见于康熙五十三年（1714）刊刻的《曝书亭集》中。这时太子已经再次被废（1712年康熙再次废黜太子）。可见康熙年间，哪怕太子再次被废，朱彝尊的“喜闻”之诗也还是可以刊刻流传的。但当《曝书亭集》在乾隆年间被收入《四库全书》时，这首“喜闻”之诗就被删除了。\NoteRef{preface}{7}这也许是我们无法从现存文集笔记中发现太多资料来了解康熙朝储位之争的一个原因。


% 原书前言第6–7页

除传统的文人资料以外，清史学者还一直在寻求域外资料中的记载。1930年代，孟森在北京大学授课所用的《清史讲义》就已经在教学生使用《朝鲜实录》来比对清朝官方历史记载。比如在涉及康熙朝废黜、复立太子事时，他已经发现了《朝鲜实录》中有一段清史完全失载的历史：康熙为了复立太子，曾经公然在朝堂上“殴曳马齐”。马齐是当时首席满洲大学士，是满人臣子的领袖。清官方正史曾记录他和康熙的舅舅佟国维数次坚持举荐皇八子为新太子，而拒绝赞同康熙复立太子的想法。《实录》中也确实能读出康熙为新太子人选与马齐等众臣几次互不相让的争执，甚至记载马齐在与康熙争执新太子人选时，曾经“拂袖而出”。但双方的矛盾是否能达到在朝堂上动手的程度？当时的康熙五十五岁，而马齐五十七岁。他们两位在这个年龄于朝堂上当着众臣之面动手，不仅有辱斯文，且完全背离了传统历史中构建出的君臣形象。起码孟森看到“殴曳马齐”的记录后表示难以相信，给出的解释是：“朝鲜忠于明”，因而“其评清帝本不甚作美辞，自难尽信”\NoteRef{preface}{8}。但“殴曳马齐”确实有戏剧冲突效果，因而后来的小说、影视剧等都选择相信，让“殴曳马齐”渐渐成为了一个进入大众认知中的有趣桥段。后来的清史学者在写到康熙朝历史时，都绕不开“殴曳马齐”这一节，用清官方史料批判此桥段之不可信。到底是朝鲜史料可信，还是清官方史料可信，很难一概而论。正巧本书使用的传教士资料将会为这一历史事件带来新的记录，也许会为康熙是否在朝堂上动手打过马齐的争论画上一个句号。


% 原书前言第7页

许多传教士在康熙一朝隶属于内务府，可以进出内廷。\NoteRef{preface}{9}传教士材料长期以来都是了解康熙朝历史的一个重要途径。早在1940年代，王锺翰在其清史成名作《清世宗夺嫡考实》一文中用到的一条关键证据，就是当时在北京的传教士马国贤留下的。\NoteRef{preface}{10}后来更多在西方做研究的学者，如徐中约（Immanuel C. Y. Hsü）、吴秀良、史景迁（Jonathan D. Spence）等进一步拓展了用传教士的记录来研究清史的方法。


% 原书前言第7–8页

从方法上来说，本书通过传教士的记录来比对正统清史资料，与前人思路并无不同。而与前人研究不同的是本书所主要依赖的资料《北京纪事》本身。有别于此前已经被广为使用的传教士信件以及他们的年度报告等，《北京纪事》有三个特点：首先，它不是回忆性的记录，而是日志型的记载。也就是说，它和现存记载康熙废立太子最主要的清官方资料《实录》拥有相同的记录模式。这样一来，两种记录的对比能让我们直观地看到二者的异同。二是由于《北京纪事》中关于天主教方面的观点不被罗马教廷认可，因而它没有像许多传教士信件一样对外刊发。《北京纪事》最近几年才陆续出版，在此前的三百年间，只有极少数的教会学者使用过它。因而对于清史研究来说，这完全是一套全新的资料。第三，也可能是最重要的一点，《北京纪事》中关于康熙废黜、复立太子期间的事件框架，以及时间、地点等都与《实录》记载一致，而不同的是事件中的人物关系，即《北京纪事》记载的几位皇子在事件中扮演的角色与《实录》有天壤之别。用最简单直白的话来说，《实录》中的好人，在《北京纪事》中是引起夺嫡之争的元凶；而《实录》中的恶人，在《北京纪事》中是最值得同情的人。在相同的事件框架下，《北京纪事》和《实录》为我们带来了两种不同但又都逻辑自洽的叙事。引人深思的是，只有一种过去曾经真实发生过。如果传教士的记录是可信的，那么之前从《实录》建立起的对于诸位皇子的认识，以及他们的历史人设，将会坍塌改变。历史前后环环相扣，相互联系。如果这段时间内《实录》树立起的皇四子、皇八子等人的历史人设崩塌了，那么我们又怎么能相信正统清史材料对接下来夺嫡争斗中各位人物的记录？而思考和回答这些问题的起点，正是验证《实录》是否存在虚构的情节。这是本书论证的核心问题，也是本书副标题中所谓“正史虚构”的根由。


% 原书前言第8–9页

本书对过去发生之事的探索主要通过比对中外材料的异同而展开。众所周知，康熙在1708年10月第一次废黜太子，1709年4月复立太子，再到1712年第二次废黜太子，此后太子位空缺，一直到1722年12月康熙去世，雍正登位。对这一段起伏的夺嫡历史，鉴于已经有众多著作进行过梳理，本书不会对此过程中的事件一一涉及。本书的主旨在于通过与太子废立相关的历史事件来观察正史中的虚构以及虚构的程度和方式。


\begin{sectionnotes}

% 原书注释：原书第387页
\BookNote{preface}{1}{AP3,513.}

% 原书注释：原书第387页
\BookNote{preface}{2}{孟森，《清世宗入承大统考实》，521—523。}

% 原书注释：原书第387页
\BookNote{preface}{3}{《大义觉迷录》版本研究及相关学术史探讨，可参见高学姣，《〈大义觉迷录〉的文本流播与记忆再现》。}

% 原书注释：原书第387页
\BookNote{preface}{4}{康熙朝储位之争是百年来的研究热点。从论证切入点大致可以分为三个方面：一是从文献所记的各种细节中，论雍正篡位；二是从皇十四子角度入手，论证雍正篡位；三是论证雍正合法继位。以下是三种论点的一些研究举例。第一种：孟森，《清世宗入承大统考实》；王锺翰，《清世宗夺嫡考实》；王锺翰，《清圣祖遗诏考辨》；杨珍，《清朝皇位继承制度》；杨珍，《关于康熙朝储位之争及雍正继位的几个问题》；许曾重，《清世宗胤禛继承皇位问题新探》；金恒源，《雍正称帝与其对手》；陈捷先，《雍正写真》；Huang Pei, Autocracy at Work: A Study of the Yung-Cheng Period 1723-1735。第二种：王锺翰，《胤祯西征纪实》；杨珍，《允禵储君地位问题研究》；庄吉发，《清世宗拘禁十四阿哥胤禵始末》；Peter C. Perdue, China Marches West: The Qing Conquest of Central Eurasia；第三种，李宪庆、白新良：《康雍之际继嗣制度的演变》；冯尔康，《雍正传》；冯尔康，《雍正继位新探》；杨启樵，《雍正篡位说驳难》；史松，《雍正研究》；吴秀良，《允禵更名与雍正继位问题再探讨》；Silas H. L. Wu, Passage to Power: K'ang-hsi and His Heir Apparent, 1661-1722; Madeleine Zelin, “Yung-cheng Reign”；常建华，《从西方文献看雍正帝继位的合法性》。}

% 原书注释：原书第388页
\BookNote{preface}{5}{根据2009年北京、台湾两地整理出版的《清代起居注册·康熙朝》。}

% 原书注释：原书第388页
\BookNote{preface}{6}{Wu, Passage to Power, 214, note 14.}

% 原书注释：原书第388页
\BookNote{preface}{7}{对于这一版本上的区别，就笔者所见最早发现者是王锺翰，《清世宗夺嫡考实》，注释26，1089。此文最早刊载于1949年的《燕京学报》。}

% 原书注释：原书第388页
\BookNote{preface}{8}{孟森，《明清史讲义》，464。}

% 原书注释：原书第388页
\BookNote{preface}{9}{详细内容参看拙作《康熙的红票：全球化中的清朝》，第四章。}

% 原书注释：原书第388页
\BookNote{preface}{10}{对于该文是“成名作”，参见王锺翰自述，《王锺翰清史论集》，2。}

\end{sectionnotes}
